% Drawing 73/01-17 -- folder 73, batch 01, archive page 18 (PDF page 19).
% Transcription: en tête, à gauche, une figure : un polygone non convexe à neuf sommets $s_0, \ldots, s_8$, de côtés $u_0, \ldots, u_7$ (le neuvième côté, de $s_8$ à $s_0$, sans lettre), avec des traits fins de $s_0$ vers $s_2$, $s_4$ et $s_7$.
% Kind: non-convex 9-gon with segments from s_0. Method 2 (schematic redraw from the facsimile). Status: ok.
% Notes: As on the page: the nine vertices computed once and every side drawn between them; the sides s_0 s_1, s_4 s_5 and s_6 s_7 heavy as on the page; thin segments from s_0 to s_4 and to s_7, and a thin one towards s_2 that stops short of it; the closing side s_8 s_0 thin (he drew it twice). Labels s_0 ... s_8 and u_0 ... u_7 where he wrote them (u_0 is written under the side s_0 s_1, beside the thin segment to s_7; s_4 and s_7 are written over earlier letters). Left out: a stray « s » left of s_4 and the doubling of the long side u_5. The transcription's note says segments go from s_0 to the other vertices: only those to s_2, s_4, s_7 (and the side to s_8) are drawn. Redrawn 2026-11-12 (issue 128): to check again.
% Coordinates: 1 unit = 1.5 mm = 1/100 of the crop width (bbox in page fractions: [0.17, 0.06, 0.43, 0.31]).
% Redrawn by Claude (claude-fable-5-2), issue #128, 2026-11-12; not checked by a human; not compiled (no TeX in the session).
% Rights: reproduces part of the work -- behind RENDER_TRANSCRIPTS, never published by itself.
\documentclass[tikz,border=3pt]{standalone}
\usepackage{amsmath,amssymb}
\usetikzlibrary{patterns,arrows.meta}
\tikzset{pencil/.style={black!45}, struck/.style={black!35}, mpt/.style={circle,draw,inner sep=1.3pt}, pt/.style={circle,fill,inner sep=0.9pt}, opt/.style={circle,draw,fill=white,inner sep=1.1pt}}
\begin{document}
\begin{tikzpicture}[x=1.5mm,y=1.5mm,line cap=round,line join=round,font=\scriptsize]
\draw[very thick] (20,-75.5) -- (39.5,-76.5);
\draw (39.5,-76.5) -- (39,-61) -- (54.5,-47.5) -- (54.7,-24.5);
\draw[very thick] (54.7,-24.5) -- (64,-10.5);
\draw (64,-10.5) -- (63,-102.5);
\draw[very thick] (63,-102.5) -- (35.5,-86.5);
\draw (35.5,-86.5) -- (35.3,-119.5);
\draw[pencil] (35.3,-119.5) -- (20,-75.5);
\draw[pencil] (20,-75.5) -- (54.7,-24.5);
\draw[pencil] (20,-75.5) -- (35.5,-86.5);
\draw[pencil] (20,-75.5) -- (36.1,-63.2);
\node[inner sep=1pt] at (19.5,-80.5) {\small $s_0$};
\node[inner sep=1pt] at (41,-80.5) {\small $s_1$};
\node[inner sep=1pt] at (43.5,-63.5) {\small $s_2$};
\node[inner sep=1pt] at (56,-53.5) {\small $s_3$};
\node[inner sep=1pt] at (59.5,-27) {\small $s_4$};
\node[inner sep=1pt] at (69,-12.5) {\small $s_5$};
\node[inner sep=1pt] at (69,-104.5) {\small $s_6$};
\node[inner sep=1pt] at (40,-95) {\small $s_7$};
\node[inner sep=1pt] at (41,-122.5) {\small $s_8$};
\node[inner sep=1pt] at (29,-82.5) {\small $u_0$};
\node[inner sep=1pt] at (43.8,-71) {\small $u_1$};
\node[inner sep=1pt] at (49.8,-57) {\small $u_2$};
\node[inner sep=1pt] at (59,-38) {\small $u_3$};
\node[inner sep=1pt] at (56.3,-15.5) {\small $u_4$};
\node[inner sep=1pt] at (70.5,-65) {\small $u_5$};
\node[inner sep=1pt] at (53.5,-93) {\small $u_6$};
\node[inner sep=1pt] at (41.5,-107.5) {\small $u_7$};
\end{tikzpicture}
\end{document}
